\documentclass[12pt,a4paper]{extarticle}
\usepackage[utf8]{inputenc}
\usepackage[T5]{fontenc}
\usepackage{geometry}
\geometry{a4paper, left=1.5cm, right=1cm, top=1.5cm, bottom=1.5cm, headheight=20pt, headsep=15pt, footskip=28pt}
\usepackage{enumitem}
\usepackage{setspace}
% Gọi package nằm trong thư mục con template
\usepackage{../template/mngocbox}

\begin{document}
\onehalfspacing

% Sử dụng ngoặc vuông [...] để thay đổi linh hoạt các thông số
\makeheadbox[headerright={Tổng ôn 10-11},chude={CHỦ ĐỀ 06},tieude={Xác suất cơ bản},dinhdang={Đại số},lop={12 },gv={Nguyễn Lê Minh Ngọc}]
\makeheading{IV}{Đề kiểm tra}
\noindent\textbf{1.} Tổ Một có 9 thành viên. Tuần tới là phiên trực nhật của tổ, nên cần phân công 4 bạn đi bê ghế của lớp cho buổi chào cờ.
\noindent a) Tổ có bao nhiêu cách phân công 4 bạn đi bê ghế?
\noindent b) Tổ có bao nhiêu cách chọn 5 bạn không phải đi bê ghế?
\vspace{0.4cm}

\noindent\textbf{2.} Chọn 4 trong số 3 học sinh nam và 5 học sinh nữ tham gia một cuộc thi.
\noindent a) Nếu chọn 2 nam và 2 nữ thì có bao nhiêu cách chọn?
\noindent b) Nếu trong số học sinh được chọn nhất thiết phải có học sinh nam $A$ và học sinh nữ $B$ thì có bao nhiêu cách chọn?
\noindent c) Nếu phải có ít nhất một trong hai học sinh $A$ và $B$ được chọn, thì có bao nhiêu cách chọn?
\noindent d) Nếu trong 4 học sinh được chọn phải có cả học sinh nam và học sinh nữ thì có bao nhiêu cách chọn?

\noindent\textbf{3.} Hãy khai triển và rút gọn biểu thức: $(1 + x)^4 + (1 - x)^4$. Sử dụng kết quả đó để tính giá trị gần đúng của biểu thức $T = 1{,}05^4 + 0{,}95^4$.

\noindent\textbf{4.} Xác định hệ số của $x^4$ trong khai triển biểu thức $(3x + 2)^5$.

\noindent\textbf{5.} Tìm hệ số của $x^2$ trong khai triển $(2x - 5)^4$.

\noindent\textbf{6} Nhân ngày khai trương siêu thị MC, các khách hàng vào siêu thị được đánh số thứ tự là các số tự nhiên liên tiếp và có thể được tặng quà (khách hàng đầu tiên được đánh số thứ tự là số 1). Cứ 4 khách vào MC thì khách thứ 4 được tặng một cái lược chải tóc, cứ 5 khách vào MC thì khách thứ 5 được tặng một cái khăn mặt, cứ 6 khách vào MC thì khách thứ 6 được tặng một hộp kem đánh răng. Sau 30 phút mở cửa, có 200 khách đầu tiên vào MC và tất cả khách vẫn ở trong MC. Chọn ngẫu nhiên 1 khách trong 200 khách đầu tiên, xác suất để chọn được khách hàng được tặng cả 3 món quà là:

\noindent\textbf{7} Có 50 tấm thẻ đánh số từ 1 đến 50. Rút ngẫu nhiên 3 thẻ. Xác suất để tổng các số ghi trên 3 thẻ chia hết cho 3 bằng:

\end{document}